\title{LongRoPE: Extending LLM Context Window Beyond 2 Million Tokens}

\begin{abstract}
Expanding the context window in large language models (LLMs) is highly sought after. Nonetheless, existing approaches are constrained to context lengths of roughly 128k tokens, primarily because of prohibitive fine-tuning expenses, limited availability of lengthy textual data, and destabilizing effects caused by unfamiliar token positions.
\end{abstract}

In this work, we present LongRoPE, a novel approach that, for the first time, expands the context window of existing pre-trained LLMs to an impressive \textbf{2048k} tokens. Remarkably, this is accomplished with no more than 1k fine-tuning steps at up to 256k training sequence lengths, while preserving the performance achieved at the standard short context window. Our method is underpinned by three main innovations: (i) we uncover and leverage two kinds of non-uniformities inherent in positional interpolation using an efficient search procedure, facilitating superior initialization for fine-tuning and supporting up to 8$\times$ context window extension without additional fine-tuning; (ii) we propose a staged extension approach where the LLM is first fine-tuned at 256k token lengths, followed by a second round of positional interpolation on this extended LLM, thereby enabling a 2048k token context window; (iii) we recalibrate LongRoPE on sequences of length 8k to restore optimal performance at short context windows. Comprehensive evaluations conducted on LLaMA2 and Mistral, spanning diverse tasks, validate the effectiveness of our approach. Models enhanced through LongRoPE maintain their original architecture with only slight adjustments to positional embeddings, and are able to reuse the majority of existing optimization techniques. The code will be released at \texttt{https://github.com/microsoft/LongRoPE}

\section{Introduction}

Although Large Language Models (LLMs) have achieved impressive results across multiple benchmarks~\cite{gpt4,llama2}, their effectiveness is hindered by constrained context window sizes, such as the 4096-token boundary in LLaMA2~\cite{llama2}. Once the input surpasses this window, the models' accuracy diminishes, since they have not encountered data beyond these positions during training. This limitation becomes particularly problematic in critical applications, including in-context learning with large quantities of examples~\cite{Huang2023FewerIM} and in the operation of LLM agents~\cite{park2023generative,madaan2023selfrefine}.

Recent studies have demonstrated that the context window of a pre-trained LLM can be increased to approximately 128k tokens by applying fine-tuning methods to longer sequences~\cite{longlora,pi,yarn,zhang2024soaring,liu2023scaling}. However, scaling the context window further presents several significant challenges. The first challenge arises from the presence of uninitialized positional indices, which can introduce anomalous outputs and lead to distributional shifts, thereby hindering the convergence of the fine-tuning process~\cite{pi}. This issue becomes particularly acute when expanding the window from 4k to over 1000k, where more than 90\% of positions are entirely new. The second challenge involves the necessity for training samples that match the desired extended context lengths. Current datasets offer very few exceptionally lengthy sequences, especially those surpassing 1000k tokens. In addition, training on such long sequences demands considerable computational power, often entailing excessive training durations and high GPU usage. Lastly, extremely long context expansions cause the attention mechanism to become increasingly diluted, distributing focus across a vast number of token positions. As a result, the model’s effectiveness on shorter contexts typically suffers~\cite{pi}.

A common strategy for addressing the initial issue involves applying interpolative techniques to RoPE positional embeddings~\cite{rope,pi}, which rescale novel position indices to fit within the range encountered during pretraining, as illustrated in Fig.\ref{fig:overview}. Position Interpolation (PI)~\cite{pi} employs linear interpolation for RoPE's rotary angles according to the ratio of extension. The NTK method~\cite{ntk,dynamicntk} introduces the concept of applying both non-uniform interpolation and extrapolation along the different RoPE dimensions. Meanwhile, YaRN~\cite{yarn} divides RoPE dimensions into three distinct groups based on frequency content and selectively assigns extrapolation, NTK, and linear interpolation procedures to each. Nonetheless, the Transformer architecture encodes positional information via \emph{complex, heterogeneous information entropy}. Current methods largely overlook this nuanced, non-uniform entropy distribution, which results in suboptimal information retention and constrains the expandability of the context window.

Section~\ref{sec:analysis} presents two main empirical insights: \textbf{(1)} Successful positional interpolation must account for two types of non-uniformities: the diverse characteristics of RoPE dimensions and the range of token positions. Minimal interpolation tends to be more beneficial for lower RoPE dimensions and tokens appearing earlier in a sequence, though the optimal degree of interpolation is contingent on the specific extension length desired. 
\textbf{(2)} Incorporating these patterns of non-uniformity into the positional interpolation process facilitates the preservation of original RoPE information, with a particular emphasis on crucial dimensions and positional indices. This approach reduces the degree of information loss attributed to positional interpolation, thereby providing a superior initialization for subsequent fine-tuning. Furthermore, it enables models to handle sequence extensions up to 8$\times$ their original length even in the absence of fine-tuning.

Building on these insights, we present \textbf{LongRoPE}, a robust approach designed to expand the LLM context window to surpass 2 \emph{million} tokens. LongRoPE is structured around three central innovations. To begin with, {LongRoPE} leverages multidimensional non-uniformities in the process of positional interpolation. It determines optimal rescale factors for the rotation angles of RoPE across different dimensions, considering the specific token positions involved. Because the space of possible rescale factors grows exponentially with increasing extension ratios, {LongRoPE} employs an evolutionary search algorithm enhanced by two optimization strategies to significantly improve the efficiency of this search. An illustration of the resulting rescaled RoPE found by the algorithm can be seen in Fig.~\ref{fig:overview}.

Subsequently, {LongRoPE} employs a resource-efficient and incremental extension method to reach a 2048k context window, circumventing the necessity of directly fine-tuning on ultra-long text sequences, which are both limited in availability and challenging to source. The approach initiates by identifying the optimal configuration at a 256k sequence length on the pre-trained LLM, which is then fine-tuned at this length. Capitalizing on our non-uniform positional interpolation mechanism, which supports up to an 8$\times$ extension without additional fine-tuning, a subsequent search is performed to determine updated RoPE rescale factors on this fine-tuned, length-extended LLM. This process enables LLaMA2 and Mistral~\cite{mistral} to support context windows up to 2048k tokens.

In order to address potential drops in performance on the initial, shorter context windows, {LongRoPE} further refines the RoPE rescale factors in the expanded LLM. Following an analogous approach to the process used when increasing from 256k to 2048k context length, the method involves reducing context windows to 4k and 8k on the 256k fine-tuned LLM. This is done by leveraging the search algorithm to minimize positional interpolation. During inference, whenever the sequence length falls below 8k, RoPE is updated using the optimized rescale factors discovered by this search.

A broad range of experiments utilizing several LLMs and a variety of long-context evaluation tasks highlight the strength of our approach. Our findings indicate that LongRoPE consistently preserves low perplexity across evaluation lengths from 4k up to 2048k tokens, attains passkey retrieval accuracy exceeding 90\%, and matches standard benchmark performance designed for the 4096 context window. Furthermore, LongRoPE is compatible with any LLM employing RoPE embeddings. We intend to make both our implementation and the LongRoPE-2048k models publicly available.

\section{Non-uniformity in Positional Interpolation}
\label{sec:analysis}

\subsection{Preliminary}

Transformer architectures depend on the provision of positional cues, typically integrated as position embeddings, to capture the sequential arrangement of input tokens. This study examines the RoPE~\cite{rope} position embedding method, a technique extensively adopted in contemporary LLMs. Specifically, for a token located at index $n$, its associated RoPE encoding is given by the simplified form:

\begin{equation}
		\fontsize{7.8}{7.8} \selectfont
	\label{eq:rope}
	\left[ cos(n\theta_0), sin(n\theta_0), cos(n\theta_1), \cdot\cdot\cdot, cos(n\theta_{d/2-1}), sin(n\theta_{d/2-1}) \right]   
\end{equation}

Here, $d$ indicates the size of the embedding dimension, 
while $n\theta_i$ corresponds to the rotary angle for a token located at position $n$, 
with $\theta_i = \theta^{-2i/d}$ defining the rotation frequency parameters. The RoPE approach adopts 10000 as the standard base value for $\theta$.

\textbf{\emph{Context window extension ratio $s$} and positional interpolation}. We define $s$ as  the ratio of extended context length $L'$ to the original length $L$: $s=\frac{L'}{L}$. 

To increase the context length from $L$ to $L'$, existing positional interpolation techniques propose adjusting the rotation frequencies $\theta_i$ by scaling them according to the context extension factor $s$. Setting $\beta = \theta^{2/d}$ and letting $\lambda$ represent the scaling factor determined by $s$, we can express these positional interpolation approaches under a unified formulation as:

\begin{equation}	
	\fontsize{7.5}{7.5} \selectfont
	\label{eq:unified_rope}
	\begin{aligned}
		\left[\cos\left(\frac{n}{\lambda(\beta)^0}\right), \sin\left(\frac{n}{\lambda(\beta)^0}\right), \cos\left(\frac{n}{\lambda(\beta)^1}\right), \cdots, \sin\left(\frac{n}{\lambda(\beta)^{d/2-1}}\right)\right]
	\end{aligned}
\end{equation}

\textbf{Linear positional interpolation (PI)}. PI~\cite{pi} employs a straightforward approach by linearly interpolating positional indices when remaining within the original sequence length used during pre-training. When extending to a sequence length scaled by a factor $s$, the rotation angles for each position across all RoPE dimensions are uniformly scaled down by $\lambda=s$. This procedure, however, compresses positional representations, resulting in denser encoding of position information and thereby reducing the model’s capacity to differentiate between tokens that are close together. As a consequence, the effectiveness of PI diminishes, particularly at larger extension ratios.

\textbf{NTK-based interpolation and extrapolation}. ~\cite{ntk,dynamicntk} approach RoPE through the lens of information encoding and leverage Neural Tangent Kernel (NTK) theory~\cite{jacot2018neural,tancik2020fourier}. To address the challenge of position crowding in PI, they propose redistributing the demands of interpolation among the dimensions of RoPE. Specifically, this approach attenuates the scaling of dimensions associated with higher frequencies while amplifying the scaling for those of lower frequencies, enabling both positional interpolation and extrapolation with the formulation $\lambda=s^{i}$. In the refined dynamic NTK~\cite{dynamicntk}, the extension ratio is dynamically modulated at each position, contingent on the sequence’s length. In contrast to PI, which relies on fine-tuning, NTK-aware strategies can lengthen context windows without the need for fine-tuning, albeit typically supporting up to a 4$\times$ extension ratio.

\textbf{YaRN}~\cite{yarn} presents a notable enhancement in positional interpolation by partitioning RoPE dimensions into three separate groups according to their frequencies, and assigning distinct interpolation methods to each group. In this approach, high-frequency dimensions are treated with extrapolation ($\lambda$=1), low-frequency dimensions employ linear interpolation (PI), and those at intermediate frequencies utilize the NTK method. The essential contribution of YaRN lies in this dimension grouping mechanism; however, it relies on manually crafted heuristics grounded in empirical analysis. As a result, this method can yield less than optimal results when applied to novel LLM architectures.

\subsection{Study on Non-uniform Positional Interpolation}

Building upon the ideas introduced in NTK and YaRN, we observe that their improvements arise from introducing non-linearity, particularly by treating the various frequencies in RoPE dimensions in a way that enables targeted interpolation and extrapolation. Despite these advances, the non-linear mechanisms employed so far are largely based on heuristics or manually crafted rules. This leads to two key questions: (1) Does the existing approach to positional interpolation represent the optimal solution? (2) Are there other forms of non-linearity that have yet to be investigated?

To address these questions, we employ evolution search (refer to Sec.~\ref{sec:method}) to identify improved non-uniform positional interpolations for LLaMA2-7B. The search process is directed by perplexity scores, utilizing five randomly selected samples from the PG19~\cite{pg19} validation set. Our empirical investigation yields the following primary insights.

\textbf{Finding 1}: \textit{RoPE dimensions exhibit substantial non-uniformities, which are not effectively handled by current positional interpolation methods.}

We optimize $\lambda$ for each RoPE dimension as specified in Eq.~\ref{eq:unified_rope}. Table~\ref{tbl:exp1} reports perplexity scores for LLaMA2-7B using various approaches on the PG19 and Proof-pile~\cite{proof-pile} test datasets, evaluated without any fine-tuning. Our proposed solution yields notable reductions in perplexity, highlighting that prevalent linear (PI) and non-uniform (Dynamic-NTK and YaRN) interpolation strategies do not achieve optimal performance. It is particularly noteworthy that YaRN performs worse than both PI and NTK on PG19, largely due to its inability to support the extended context window length in language models without fine-tuning. For instance, when tested with a context size of 8k, the perplexity for YaRN increases sharply after reaching 7k.

In our investigation, the rescaling coefficients $\lambda$ in Eq.~\ref{eq:unified_rope} exhibit variability, contrasting with the constant scaling factor $s$ used in the calculations for PI, NTK, and the group-based approach in YaRN. Utilizing these non-uniform scaling factors yields substantial enhancements in LLaMA2's language modeling capabilities—as indicated by perplexity scores—over 8k and 16k context windows, all achieved without the need for additional fine-tuning. This improvement is attributed to the positional embeddings' enhanced ability to maintain the integrity of the original RoPE structure, particularly with respect to crucial dimensions, which aids the LLM in more effectively discriminating between similar token positions.

\textbf{Finding 2}: \textit{RoPE for the initial tokens in the input sequence should be extrapolated with less interpolation.} 

We propose that, for the first $\hat{n}$ tokens in the input, reducing interpolation in their RoPE representations may be advantageous. The rationale is that these tokens are subject to higher attention weights, thus playing a pivotal role in attention mechanisms, as evidenced by findings in Streaming LLM~\cite{streamingllm} and LM-Infinite~\cite{han2023lminfinite}. To test this proposition, we conduct experiments extending the context window to 8k and 16k tokens via PI and NTK methods, while varying the count of initial tokens $\hat n$ (with values 0, 2, ..., 256) left unaltered by interpolation. For the case of $\hat n$ = 0, the standard PI and NTK configurations are used. The results, as summarized in Table~\ref{tbl:tokenposition}, yield two main insights: \textbf{(1)} exempting the initial tokens from position interpolation yields noticeable performance gains; \textbf{(2)} the most beneficial value of $\hat n$ varies with the intended length of context window expansion.

\textbf{Finding 3}: \textit{Non-uniform positional interpolation effectively extends LLM context window in both fine-tuning and non-fine-tuning settings.} 

Our experiments demonstrate that applying our optimized non-uniform position interpolation notably enhances out-of-the-box extension capabilities at 8k and 16k context lengths. However, scaling to even longer contexts necessitates additional fine-tuning. Therefore, we fine-tune LLaMA2-7B integrated with our customized RoPE approach for a 64k context window (refer to Appendix for detailed configuration). As presented in Table~\ref{tbl:finetune}, our strategy achieves substantial gains compared to both PI and YaRN, observable both prior to and following the fine-tuning phase for LLaMA2-7B. This improvement stems from leveraging non-uniform positional interpolation, which effectively preserves information and furnishes more optimal initial parameters for fine-tuning.

\textbf{Summary}. Our research identifies two non-uniform aspects: heterogeneity in RoPE dimensions and token positions. Strategic exploitation of these non-uniformities within positional interpolation frameworks yields significant advances in large language model context window extension.

\section{LongRoPE}

\label{sec:method}
Building upon these insights, we propose LongRoPE, which initially employs an optimized search strategy to leverage both types of non-uniformities, and subsequently applies this approach to expand the LLM context window past 2 million tokens.

\subsection{Problem Formulation}

The presence of these two non-uniformities dramatically expands the solution space and adds significant challenges to the optimization process. To tackle this, we cast the multidimensional non-uniform position interpolation optimization as a search-based problem.

Given an LLM with a context window of length $L'$ and an input consisting of long documents $\mathbf{X}$, such that each $\mathbf{x} \in \mathbf{X}$ exceeds $L'$ tokens, let us represent the original rotary angle in the $i^{th}$ RoPE embedding dimension at position $n$ by $\frac{n}{\beta^i}$. The resulting optimization task can be stated as follows:

\begin{equation}
\begin{gathered}
		\label{eq:problem}

\underset {\mathbf{x}\in\mathbf{X}; \, |\mathbf{x}|\ge L'}{ \operatorname {arg\,min} } \,  \mathcal{L}\, \left( \text{LLM} (\text{RoPE}, \mathbf{X})\right),\quad \text{where}
\\
\scalebox{0.93}{
$\underset { \begin{subarray}{l} i = 0,\ldots,\frac{d}{2}-1; \\ n \in [0,|\mathbf{x}|); \end{subarray} }{ \text{RoPE}_{(n)} } = \left[ \ldots, \cos \left( \mathbb{I}(\hat{\lambda}_i, \hat{n}) \frac{n}{\beta^i} \right), \sin \left( \mathbb{I}(\hat{\lambda}_i, \hat{n}) \frac{n}{\beta^i} \right), \ldots \right] $}
\\
\text{with} \quad \mathbb{I}(\hat{\lambda}_i, \hat{n}) = \begin{cases}
1 &\text{if}\;\; n < \hat{n} \\
\frac{1}{\lambda_i} &\text{if}\;\; n \geq \hat{n}
\end{cases}
	\end{gathered}
\end{equation}

We introduce a collection of scaling factors, $\mathbb{I}(\hat{\lambda}_{i},\hat{n})$, designed to address two types of non-uniformities. Here, $\hat{\lambda_i}$ corresponds to non-uniformity across RoPE dimensions, while $\hat{n}$ captures positional non-uniformity among tokens. The function $\mathbb{I}(\hat{\lambda}_{i},\hat{n})$ specifically rescales the rotation angle associated with the $i^{th}$ RoPE dimension by incorporating $\hat\lambda_{i}$ as the scaling parameter and $\hat{n}$ as the positional threshold. For the first $\hat{n}-1$ token positions, the rescaling factor $\hat\lambda_{i}$ is inactive, resulting in use of the original RoPE rotary angle $\frac{n}{\beta^i}$. When the token position satisfies $n\ge\hat{n}$, the rescale factor begins to influence the rotation.

For a desired context window of size $L'$, the goal is to determine the optimal set of rescale factors, denoted as ($\mathbb{I}(\hat{\lambda}_{0},\hat{n})$, $\mathbb{I}(\hat{\lambda}_{1},\hat{n})$, ... $\mathbb{I}(\hat{\lambda}_{i},\hat{n})$, ...), corresponding to each RoPE dimension from $1$ to $d$. By applying these rescale factors to the $\text{RoPE}$ mechanism within the target $\text{LLM}$, we aim to reduce the next token prediction loss, $\mathcal{L}$ (i.e., perplexity), to its minimum for input samples $\mathbf{X}$ of token length $L'$.

\subsection{Searching the Non-uniform Position Interpolation}

To address the challenge outlined in Eq.~\ref{eq:problem}, we propose an intuitive yet robust approach that identifies the optimal RoPE rescale factors, thereby enabling comprehensive utilization of the inherent multidimensional non-uniformities within position embeddings.

\textbf{Search space}. Our methodology establishes a broad search space designed to accommodate the two primary types of non-uniformity. Table~\ref{tbl:searchspace} delineates the structure of this search space. Concretely, we enable the optimization process to determine a distinct rescale factor for every dimension present in the RoPE embedding. For the sake of streamlined search space construction, we opt to search over $\lambda_i$ and $\hat{n}$ instead of directly searching the indicator function $\mathbb{I}(\hat{\lambda}_{i},\hat{n})$, wherein $\hat{\lambda}_{i}=1/\lambda_i$. According to the specifications in Table~\ref{tbl:searchspace}, the search range for each $\lambda_i$ extends from a minimum of 1.0 (representing straightforward extrapolation) up to a maximum of $s\times1.25$ (corresponding to a broader interpolation range than PI), with increments of 0.01, where $s$ denotes the ratio by which the context window is to be enlarged.

The parameter $\hat{n}$ specifies how many of the initial token positions maintain their original encoding without employing position interpolation, thus preserving the standard RoPE embedding. In practice, we allow $\hat{n}$ to be selected from the set \{0, 1, 2, 4, 8, 12, 16, 20, 24, 28, 32, 64, 128, 256\}. If $\hat{n}=0$, this indicates that rescale factors identified through the search are applied to every token position.

\textbf{Evolution-based search}. The expansive search space presented in Table~\ref{tbl:searchspace} encompasses a wide array of positional interpolation configurations, making systematic exploration notably challenging. For instance, when using an extension ratio of $s=4\times$, there exist $400^{128/2}\times14=4\times10^{167}$ possible combinations. The search space increases exponentially with larger extension factors. To manage this complexity, we employ an evolution search strategy~\cite{guo2020single} and propose two optimization methods that substantially enhance the efficiency of the search process. The comprehensive steps of the search strategy are provided in Algorithm~\ref{alg:search}.

\textit{Optimized initial population generation}. In place of assigning all $P$ rescale factors at random, our approach begins by explicitly including the three RoPE rescale factors associated with PI, NTK, and YaRN as part of the initial population. For the other $P$-3 members, these three rescale factors are subjected to random mutations, each with probability $p$.

\textit{Monotonically non-decreasing constraint}. Upon initializing the population, we evaluate the LLM perplexity associated with each candidate. This involves applying the respective RoPE rescale parameters to the target LLM and determining the perplexity over input $\mathbf{X}$. The highest-performing $k$ individuals are chosen to serve as parents in the evolutionary process. Nonetheless, due to the extensive nature of the search space, unrestrained mutation and crossover may frequently yield suboptimal candidates, resulting in superfluous perplexity evaluations. Such inefficiency is exacerbated for large values of $L'$, since each perplexity measurement demands substantial inference time.

To tackle this issue, we enforce a constraint ensuring that the sequence of rescaled RoPE factors is non-decreasing, such that $\lambda_i \leq \lambda_{i+1}$. Only those RoPE configurations meeting this monotonicity requirement are considered for perplexity computation in the LLM, which markedly cuts down the search space. Concretely, we stipulate that each $\lambda_i$ must grow monotonically across increasing RoPE dimensions (i.e., for $i=0,...,63$). This form of monotonicity across dimensions is informed by NTK theory~\cite{jacot2018neural,tancik2020fourier,ntk}, which posits that lower-dimensional components associated with higher frequencies benefit from minimal interpolation (thus, a smaller $\lambda_i$), while higher-dimensional components, corresponding to lower frequencies, tolerate greater degrees of interpolation (therefore, a larger $\lambda_i$).

\begin{algorithm}[t]
	\small
	\caption{The search algorithm}
	\textbf{Input:} target LLM, input samples $\mathbf{X}$,   population size $P$, mutation size $N_1$, crossover size $N_2$, max iterations $\mathcal{T}$, mutate probability $p$\\

	\begin{algorithmic}[1]
		\label{alg:search}
		\STATE Initialize Top-k as $\phi$;	
		\STATE Set $\text{P}_0$ by calling \textit{Initialize\_population\_with\_optimization} ($P$, $\mathbf{X}$, $p$); 
		\FOR{$i=1$ to $\mathcal{T}$}
		\STATE Run \textit{Compute\_perplexity} (LLM, $\text{P}_{i-1}$, $\mathbf{X}$);
		\STATE  Update Top-k via \textit{Update\_Topk} (Top-k, $\text{P}_{i-1}$);
		\STATE Obtain $\text{P}_{mutation}$ using \textit{Mutation\_with\_mono\_constraint} (Top-k, $N_1$, $p$); 
		\STATE Derive $\text{P}_{crossover}$ through \textit{Crossover\_with\_mono\_constraint} (Top-k, $N_2$);
		\STATE Form the next population: $\text{P}_i = \text{P}_{mutation} \cup \text{P}_{crossover} \cup$ Top-k;
		\ENDFOR
		\STATE Output the element in Top-k with the minimum perplexity;
	\end{algorithmic}
\end{algorithm}

\textbf{8$\times$ extension without fine-tuning}. Through evolutionary search, our approach successfully discovers non-uniform RoPE rescale factors that strategically maintain critical dimensions and positions, thereby reducing information degradation caused by interpolation. As shown in Fig.\ref{fig:non-ft}, this strategy enables the LLaMA2 model to expand its context window from 4k to 32k tokens without any additional fine-tuning. By comparison, prior methods like PI, as well as non-uniform NTK and YaRN, experience sharp increases in perplexity once the context window is doubled.

\subsection{Extending LLM Context Window to 2048K}
\label{sec:extendto2048k}

\textbf{Progressive extension to 2048k}. In this section, we present our approach for expanding the context window of pre-trained LLMs from the typical 4k tokens up to over 2048k. Our non-uniform positional interpolation technique is shown to allow an 8$\times$ increase in context length without the need for fine-tuning. However, for more significant extensions (e.g., up to 512$\times$), fine-tuning becomes necessary. One possible strategy involves determining appropriate RoPE rescaling factors tailored to the 2048k target and subsequently fine-tuning the model. Nevertheless, this option is hindered by the extremely high computational demands of training. In addition, our observations suggest that achieving effective fine-tuning becomes increasingly difficult as the extension ratio grows larger (refer to Appendix for further details).

Fortunately, LongRoPE demonstrates strong performance on both the base model and the extended LLM after fine-tuning. As a result, we propose a streamlined, incremental approach capable of reaching the desired 2048k length in only 1k fine-tuning steps, all within a maximum training sequence length of 256k.

$\Diamond$ \textit{LongRoPE search for expanding pre-trained LLMs to 256k}. Using LLaMA2 as the representative case, we systematically explore configurations for context window lengths of 128k and 256k. At this phase, the model undergoes expansion by factors of 32$\times$ and 64$\times$, respectively.

$\Diamond$ \textit{Fine-tuning to 256k}. In this phase, the pre-trained LLM undergoes additional training to reach a 256k context window. Initially, LLaMA2 is fine-tuned for 400 steps with RoPE rescaled factors configured for 128k. After this stage, we update the RoPE rescaled factors to target 256k on the resulting checkpoint and proceed with a further 600 fine-tuning steps. This two-stage approach is more effective than immediately fine-tuning to the 256k context length.

$\Diamond$ \textit{Scaling the context window of a fine-tuned extended LLM to 2048k via LongRoPE search}. In the final phase, we conduct an additional search procedure over the LLM that was previously fine-tuned with a 256k sequence length. This approach enables us to expand the context window dramatically to 2048k tokens, achieving this substantial increase without necessitating extra fine-tuning steps. As a result, the overall context extension reaches $512\times$.

\textbf{Performance restoration for shorter context windows}. Expanding the context window to an extensive 2048k length leads to decreased performance when evaluating within the initial context span. This phenomenon is a recognized limitation of positional interpolation~\cite{pi}, where embedding positions in the higher-dimensional space for the initial context are compressed into a restricted area, thereby impairing the language model’s capabilities. When the extension ratio reaches 512$\times$, the positions corresponding to the original 4k context window are especially densely packed.

To address this issue, we conduct an additional evolution search specifically for the extended LLM, refining the RoPE rescale factors tailored to shorter context windows such as 4k and 8k. For these shorter lengths, we limit the upper bound for the searched $\lambda$, since positional interpolation needs are less pronounced. When the model is used for inference, it adapts the relevant RoPE rescale factors on-the-fly.

\section{LongRoPE}

\label{sec:method}
Building on these insights, we propose LongRoPE, a method that incorporates an optimized search algorithm specifically designed to leverage both types of non-uniformities. Using this approach, we are able to expand the LLM context window to lengths exceeding 2 million tokens.

\subsection{Problem Formulation}

The presence of these two forms of non-uniformity considerably expands the solution landscape and adds layers of difficulty to the optimization process. To tackle this challenge, we recast the multidimensional non-uniform position interpolation task as a search-based optimization problem.

Consider a large language model designed for a context window length of $L'$, and a set of extended input texts $\mathbf{X}$, where each sequence $\mathbf{x}\in\mathbf{X}$ has more tokens than $L'$. We express the standard rotary angle for the $i^{th}$ component in the RoPE embedding at token index $n$ as $\frac{n}{\beta^i}$. The corresponding optimization task can be expressed in the following way:

\begin{equation}
\begin{gathered}
		\label{eq:problem}

\underset{\mathbf{x}\in\mathbf{X};\,|\mathbf{x}|\geq L'}{\operatorname{arg\,min}}\, \mathcal{L}\, \left( \text{LLM} (\text{RoPE}, \mathbf{X})\right), \text{where }
\\
\scalebox{0.93}{
$\underset{\begin{subarray}{l} i=0,\dots, \frac{d}{2}-1;\\ n\in[0, |\mathbf{x}|);\end{subarray}}{\text{RoPE}_{(n)}} = {\left[\dots, \cos\left(\mathbb{I}(\hat{\lambda}_{i}, \hat{n})\cdot\frac{n}{\beta^i}\right), \sin\left(\mathbb{I}(\hat{\lambda}_{i}, \hat{n})\cdot\frac{n}{\beta^i}\right), \dots\right]}$}
\\
\text{with } \mathbb{I}(\hat{\lambda}_{i},\hat{n}) =
\begin{cases}
1 & \text{if } n < \hat{n} \\
\frac{1}{\lambda_{i}} & \text{if } n \geq \hat{n}
\end{cases}
\end{gathered}
\end{equation}

In this approach, we define a group of rescaling coefficients, $\mathbb{I}(\hat{\lambda}_{i},\hat{n})$, intended to address two distinct kinds of non-uniformities. Here, $\hat{\lambda}_{i}$ characterizes the non-uniform scaling within the RoPE dimensions, while $\hat{n}$ specifies a positional threshold for tokens. The function $\mathbb{I}(\hat{\lambda}_{i},\hat{n})$ is applied to adjust the rotational angle specifically for the $i^{th}$ RoPE dimension: $\hat{\lambda}_{i}$ serves as the rescaling parameter, and $\hat{n}$ determines when this parameter becomes active. For the first $\hat{n}-1$ token positions, the rescaling $\hat{\lambda}_{i}$ does not alter the computation, so the conventional RoPE rotation angle $\frac{n}{\beta^i}$ remains unchanged. Beyond this, i.e., for tokens with positions $n \geq \hat{n}$, the rescaling factor is introduced.

Assuming a desired context window length of $L'$, the task is to determine the best set of rescale factors ($\mathbb{I}(\hat{\lambda}_{0},\hat{n})$, $\mathbb{I}(\hat{\lambda}_{1},\hat{n})$, ...$\mathbb{I}(\hat{\lambda}_{i},\hat{n})$...) across the $1^{st}$ through $d^{th}$ RoPE dimensions. By applying these rescale factors to the target $\text{LLM}$'s $\text{RoPE}$ component, we aim for the lowest possible next token prediction loss, $\mathcal{L}$ (which corresponds to perplexity), on input sequences $\mathbf{X}$ of length $L'$.

\subsection{Searching the Non-uniform Position Interpolation}

To address the issue outlined in Eq.~\ref{eq:problem}, we propose a straightforward yet powerful approach that involves identifying the most suitable RoPE rescale factors, thereby harnessing the inherent multidimensional asymmetries present in position embeddings.

\textbf{Search space}. We construct an extensive search space that accounts for both types of non-uniformity. The details of this search space are summarized in Table~\ref{tbl:searchspace}. In particular, we permit the selection of distinct rescale factors for each individual dimension within the RoPE embedding. For the sake of simplifying the search space formulation, our method focuses on exploring $\lambda_i$ and $\hat{n}$ instead of directly optimizing over $\mathbb{I}(\hat{\lambda}_{i},\hat{n})$, where  $\hat{\lambda}_{i}=1/\lambda_i$. As indicated in Table~\ref{tbl:searchspace}, the range for $\lambda_i$ starts from a lower bound of 1.0 (representing straightforward extrapolation), extending up to $s\times1.25$ (which allows for even greater interpolation than PI), with increments of 0.01. Here, $s$ denotes the desired context window extension ratio.

The parameter $\hat{n}$ determines how many of the earliest token positions are kept using the unaltered RoPE embeddings, meaning no position interpolation is applied to them. In our experiments, $\hat{n}$ is selected from the set \{0, 1, 2, 4, 8, 12, 16, 20, 24, 28, 32, 64, 128, 256\}. If $\hat{n}=0$, then every token position employs the optimized rescale factors.

\textbf{Evolution-based search}. The search space described in Table~\ref{tbl:searchspace} encompasses a vast array of possible positional interpolation configurations, making thorough and efficient exploration difficult. For instance, when the extension ratio is set to $s=4\times$, the total number of possible selections reaches $400^{128/2}\times14$ = 4$\times10^{167}$. The complexity of the search increases exponentially as the extension ratio grows. To navigate this extensive space, we employ evolution search~\cite{guo2020single} and incorporate two key optimization strategies that significantly improve the efficiency of the search process. The complete search method is detailed in Algorithm~\ref{alg:search}.

\textit{Optimized initial population generation}. Rather than relying solely on random selection for the $P$ rescale factors constituting the initial population, we explicitly include the three RoPE rescale factors associated with PI, NTK, and YaRN as dedicated individuals. To populate the remaining $P$-3 individuals, we introduce stochastic variations by randomly mutating these three rescale factors with a mutation probability of $p$.

\textit{Monotonically non-decreasing constraint}. Following the creation of the initial population, we evaluate the LLM perplexity for every candidate. To do this, we assign the appropriate RoPE rescale factors within the target LLM and determine the perplexity for the input $\mathbf{X}$. The best $k$ individuals are then selected as parents for the subsequent evolutionary steps. Due to the expansive nature of the search space, indiscriminate mutation and crossover operations may result in the assessment of suboptimal candidates, thereby incurring superfluous perplexity evaluations. This inefficiency is exacerbated when $L'$ is large, since each perplexity calculation requires substantial inference time.

To address this issue, we enforce a monotonic non-decreasing constraint on the set of sampled RoPE rescaling factors: $\lambda_i\le \lambda_{i+1}$. Only those RoPE configurations that fulfill this restriction are used for LLM perplexity evaluations, leading to a substantial reduction in the cost of the search. More concretely, we mandate that $\lambda_i$ increases monotonically as a function of RoPE dimension ($i = 0, \ldots, 63$). This condition is motivated by NTK theory~\cite{jacot2018neural,tancik2020fourier,ntk}, which indicates that lower dimensions—corresponding to higher frequencies—require less interpolation (i.e., a smaller value of $\lambda_i$), whereas higher dimensions—with lower frequencies—can accommodate greater degrees of interpolation (i.e., a larger $\lambda_i$).

\begin{algorithm}[t]
	\small
	\caption{The search algorithm}
	\textbf{Input:} target LLM, input samples $\mathbf{X}$, population size $P$, mutation size $N_1$, crossover size $N_2$, maximum iterations $\mathcal{T}$, mutation probability $p$ \\

	\begin{algorithmic}[1]
		\label{alg:search}
		\STATE Top-k $\leftarrow \phi$;	
		\STATE $\text{P}_0 \leftarrow$\textit{Initialize\_population\_with\_optimization}($P$, $\mathbf{X}$, $p$); 
		\FOR{$i = 1$ \textbf{to} $\mathcal{T}$}
		\STATE \textit{Compute\_perplexity}(LLM, $\text{P}_{i-1}$, $\mathbf{X}$);
		\STATE Top-k $\leftarrow$ \textit{Update\_Topk}(Top-k, $\text{P}_{i-1}$);
		\STATE $\text{P}_{mutation} \leftarrow$ \textit{Mutation\_with\_mono\_constraint}(Top-k, $N_1$, $p$); 
		\STATE $\text{P}_{crossover} \leftarrow$ \textit{Crossover\_with\_mono\_constraint}(Top-k, $N_2$);
		\STATE $\text{P}_i \leftarrow \text{P}_{mutation} \cup \text{P}_{crossover} \cup$ Top-k;
		\ENDFOR
		\STATE \textbf{Return} the candidate with the minimal perplexity from Top-k;
	\end{algorithmic}
\end{algorithm}

\textbf{8$\times$ extension without fine-tuning}. Through evolutionary search, our approach efficiently determines non-uniform RoPE rescale factors, selectively retaining critical dimensions and locations to reduce information loss caused by interpolation. As shown in Fig.\ref{fig:non-ft}, this strategy enables the expansion of LLaMA2's context window from 4k up to 32k without requiring any fine-tuning. By comparison, alternative techniques such as PI, as well as non-uniform NTK and YaRN, experience a significant increase in perplexity following a 2$\times$ extension.

\subsection{Extending LLM Context Window to 2048K}
\label{sec:extendto2048k}

\textbf{Progressive extension to 2048k}. In this section, we present our approach for scaling the context window of pre-trained LLMs beyond the standard 4k tokens, reaching up to over 2048k. As shown earlier, the application of our non-uniform positional interpolation method enables an 8$\times$ increase in context length without necessitating additional fine-tuning. When aiming for substantially larger extensions (for instance, 512$\times$), fine-tuning becomes indispensable. A commonly used strategy is to identify suitable RoPE rescaling parameters for the targeted 2048k context length, followed by fine-tuning. However, this strategy is constrained by the considerable computational resources required. Furthermore, our empirical observations indicate significant difficulty in effectively fine-tuning LLMs for such large-scale extension ratios (refer to Appendix).

Luckily, LongRoPE proves to be successful with both the base and the fine-tuned, length-extended LLM. Thus, we propose a streamlined, stepwise approach that reaches the desired 2048k target using merely 1k fine-tuning iterations, all conducted within a maximum training length of 256k.

$\Diamond$ \textit{Investigating the expansion of pre-trained LLMs to 256k via LongRoPE search}. Using LLaMA2 as a case study, we perform a search targeting context window sizes of 128k and 256k tokens. At these respective settings, the extension ratio reaches 32$\times$ and 64$\times$.

$\Diamond$ \textit{Fine-tuning to 256k}. In the subsequent phase, the pre-trained LLM is adapted to reach a context window of 256k. The process begins with a 400-step fine-tuning of LLaMA2, utilizing RoPE rescaled factors adjusted for 128k. After completing this stage, the RoPE rescaled factors are updated to support 256k, and a further 600 fine-tuning steps are carried out on the obtained checkpoint. Compared to direct fine-tuning for 256k, this approach demonstrates improved efficiency.

$\Diamond$ \textit{Scaling fine-tuned extended LLM to 2048k via LongRoPE search}. In the final step, we conduct an additional search process using the LLM that has already been fine-tuned for a 256k context length. Remarkably, this approach allows the context window to be expanded to 2048k, achieving this increase without necessitating additional fine-tuning. This process yields a context window extension factor of $512\times$.

\textbf{Recovery within a reduced context window}. When scaling up the context window to a very large 2048k size, a decline in model performance is observed for inputs that fit within the model’s initial context length. This outcome is a documented limitation of positional interpolation~\cite{pi}: the process compresses positional embeddings within the original window into a tighter space in the expanded configuration, which impairs the language model’s effectiveness. In the case of a 512$\times$ extension, the positional indices of the initial 4k context window are densely clustered, which further exacerbates the issue.

To address this issue, we conduct an additional evolution search on the extended LLM, specifically targeting adjustments of the RoPE rescale factors at shorter context lengths (such as 4k and 8k). Since less positional interpolation is necessary at these lengths, we lower the maximum permitted $\lambda$ values in the search process. The LLM subsequently modifies the relevant RoPE rescale factors on-the-fly during inference.

 \section{Experiments}

\label{sec:exp}
\subsection{Setup}
\textbf{Evaluation Tasks and Models}. LongRoPE is implemented on LLaMA2-7B and Mistral-7B models, with their capabilities assessed across three main dimensions: (1) the perplexity of large language models (LLMs) with expanded context lengths when processing lengthy documents; (2) a Passkey retrieval evaluation, designed to quantify the model's proficiency in locating a specified passkey amidst extensive unrelated text; and (3) performance on conventional LLM benchmarks limited to a 4096-token context window.

\textbf{Fine-tuning}. For the LLaMA2 model, we adopt a linear learning rate decay schedule starting from 2e-5, utilizing a total batch size of 32 across all devices. The initial fine-tuning involves 400 optimization steps on the Redpajama~\cite{redpajama} dataset, partitioned into 128k-length segments and framed with the BOS and EOS markers. After this phase, continuing from the resulting checkpoint, we extend training for an additional 600 steps to expand the context window to 256k. Fine-tuning with a 128k context window employs 8 A100 GPUs leveraging the distributed framework~\cite{cube}, whereas extending to 256k necessitates 16 A100 GPUs. For the Mistral model, the procedure uses a constant learning rate of 1e-6 with a global batch size set to 64. The protocol outlined in YaRN~\cite{yarn} guides both 128k and 256k context size variants, each fine-tuned with 400 steps on the Together Computer's Long-Data Collections~\cite{mistraldata}, adopting a 16k token sequence length. Model training for Mistral is carried out using 4 A100 GPUs.

\textbf{Search}. For a target window size up to 256k, the following hyperparameters are adopted: $P$=64, $N_1$=$N_2$=16, $p$=0.3, and $\mathcal{T}$=40, with the top 32 individuals selected for mutation and crossover in each cycle. Perplexity scores are obtained from five randomly chosen samples of the PG19 validation set, each sample having at least the specified target context length. When handling window sizes greater than 512k, the population size as well as the number of mutation and crossover operations are reduced by half. In these cases, perplexity is evaluated using three randomly selected validation samples from Pile-Books3~\cite{pile}.

\textbf{Baselines}. In order to attain a 2048k context window, we conducted fine-tuning on models utilizing 128k and 256k context lengths. This procedure results in LongRoPE-2048k (ft=128k) for LLaMA2 and LongRoPE-2048k (ft=256k) for Mistral. We benchmark these four models against leading context window extension baselines by evaluating them alongside publicly available LLMs that were fine-tuned using positional interpolation techniques, such as PI, NTK, and YaRN. Among the evaluated baselines are Together-32k~\cite{together}, Code LLaMA~\cite{codellama}, LongLoRA-full-FT-100k~\cite{longlora}, YaRN-LLaMA, and YaRN-Mistral~\cite{yarn}.

\subsection{Main Results}

\label{sec:mainresult}
\textbf{Long sequence language modeling up to 256k tokens}. To start, we assess performance against leading extended LLMs on tasks restricted to a 256k token evaluation window. For a robust demonstration of the model’s ability to generalize, two test datasets are utilized: Proof-pile~\cite{pg19} and the test portion of PG19~\cite{pile}. Perplexity is measured at different context sizes via a sliding window approach with a window size of 256. Specifically, all 100 test documents from PG19 are included in the evaluation. For Proof-pile, consistent with the protocol of YaRN~\cite{yarn}, we randomly draw 10 samples, ensuring each sample has a minimum length of 128k tokens.

Table~\ref{tbl:proof-pile} and Table~\ref{tbl:pg19} present a perplexity comparison between LLaMA2 and Mistral models extended with various interpolation techniques, evaluated on Proof-pile and PG19, respectively. Two primary findings emerge: \textbf{(1)} our extended models consistently achieve lower perplexity as the evaluation sequence length increases from 4k to 256k, indicating improved capacity to utilize extended context. \textbf{(2)} Despite being evaluated with a context window up to 16$\times$ longer—a scenario that usually impairs performance at more moderate lengths—our LongRoPE-2048k variants achieve superior perplexity compared to state-of-the-art baselines within the 256k range.

\textbf{Long sequence language modeling beyond 2000k}. To assess performance on ultra-long documents, we utilize the Books3~\cite{pile} dataset. For practical evaluation, we randomly choose 20 books, each with a length greater than 2048k, and apply a 256k sliding window during analysis.

As indicated in Table~\ref{tbl:books3}, LongRoPE effectively extends the context window of LLaMA2-7B and Mistral-7B models to 2048k tokens and also attains perplexity values that are either on par with or better than baseline methods for shorter context lengths between 8k and 128k. We further note significant disparities in performance when comparing LLaMA2 at 2048k to Mistral. Mistral demonstrates strong results relative to baselines at short ranges, yet its perplexity rises above 7 for context lengths exceeding 256k. The perplexity metrics for LLaMA2 generally meet expectations, displaying steady improvement as the context length grows, with only modest increases at 1024k and 2048k. Additionally, for LLaMA2, fine-tuning LongRoPE-2048k with a 256k context window yields better results than with 128k, attributed to a lower secondary extension ratio (specifically, 8$\times$ compared to 16$\times$). By contrast, Mistral achieves superior performance when fine-tuned with a 128k window. This can primarily be attributed to our adoption of YaRN's procedure in which both 128k and 256k Mistral fine-tuning employ a 16k training length, a choice that impacts Mistral’s capacity to further lengthen its context window post fine-tuning.

\textbf{Passkey retrieval}. Next, we examine how well models handle retrieval tasks over varying context window sizes. Utilizing a synthetic benchmark inspired by ~\cite{passkey}, we evaluate model performance in extracting a concealed passkey—a five-digit numeral—embedded somewhere within an extended document. The prompt format utilized for this task is provided in the appendix. To assess performance, we conduct 10 independent runs, each time randomly positioning the passkey at a location uniformly sampled from the full evaluation context range.

Fig.~\ref{fig:passkey} presents a comparison of retrieval accuracy between our models and baseline methods. The performance of existing approaches declines precipitously to zero when sequence lengths exceed 128k tokens. Conversely, our LongRoPE-LLaMA2-2048k (ft=256k), even when challenged with passkey retrieval from millions of tokens, consistently delivers high retrieval accuracy ($\ge$90\%) across the 4k to 2048k range. Similarly, LongRoPE-Mistral-2048k (ft=128k) sustains perfect (100\%) accuracy up to 1800k tokens and then experiences a decrease to 60\% at 2048k. This drop is consistent with the trends seen in Table~\ref{tbl:books3}, which indicates a modest rise in perplexity at the 2048k token mark.

\textbf{Evaluation on standard benchmarks within the initial context window}. We assess LongRoPE-2048k models on the original context window utilizing the Hugging Face Open LLM Leaderboard~\cite{huggingfaceboard} across both zero-shot and few-shot paradigms. Specifically, our evaluation involves 25-shot ARC-Challenge~\cite{allenai:arc}, 10-shot HellaSwag~\cite{zellers2019hellaswag}, 5-shot MMLU~\cite{mmlu}, as well as 0-shot TruthfulQA~\cite{lin2021truthfulqa}.

As demonstrated in Table~\ref{tbl:commonsense}, our models perform on par with existing benchmarks tailored for shorter context lengths, and surpass the original Mistral on TruthfulQA by a margin of +0.5\%. While LongRoPE-LLaMA2-2048k, which underwent fine-tuning at 256k, experiences a marginally greater decline in performance, its results for the majority of tasks still fall within acceptable limits.

\subsection{Ablation Results}

\textbf{Effectiveness of the second positional interpolation}. Within our progressive extension framework, we employ our search algorithm to perform a second round of non-uniform positional interpolation on large language models that have already been fine-tuned and extended. To assess the benefits of this approach, we carry out experiments on our fine-tuned LLaMA2-256k model, extending it further to context lengths of 512k, 1024k, and 2048k via PI and YaRN methods. As indicated in Table~\ref{tbl:secondsearch}, our approach utilizing non-uniform positional interpolation maintains stable perplexity scores across these greater context sizes. Conversely, perplexity values for models extended through PI and YaRN show a sharp escalation as the extension ratio increases.

\textbf{Effectiveness of recovery at shorter context lengths.} In order to address the degradation in performance observed with shorter context lengths, we modify the RoPE scaling parameters for LongRoPE-2048k using our optimization method. More precisely, we limit the upper bounds for scale factors within the search to favor reduced interpolation when operating at 4k and 8k input lengths. Table~\ref{tbl:shortperformance} presents perplexity metrics for LongRoPE-LLaMA2-2048k evaluated on Proof-pile at 4k and 8k tokens, together with the corresponding average LLM benchmark scores. These results highlight notable gains in performance when processing short context sequences.

\textbf{Analysis on the two forms of non-uniformities}. In this section, we investigate the distinct impacts of the two types of non-uniformities through ablation studies, evaluating the contribution of each component to overall model performance. Specifically, we conduct two sets of experiments: (i) expanding LLaMA2-7B to short sequence lengths of 16k and 32k using several strategies—PI, optimizing the RoPE dimension alone, and jointly optimizing both non-uniformities; (ii) scaling our 256k-finetuned LLaMA2 up to 2048k, employing the same methodologies. Perplexity is measured prior to any further fine-tuning. According to Table~\ref{tbl:searchdimension}, modifying non-uniformity within the RoPE dimension yields a notable perplexity decrease in comparison to the linear interpolation approach used by PI. While adjusting non-uniformity along token positions clearly enhances outcomes for 16k and 32k sequences, this effect diminishes at the much larger 2048k context, which may be attributable to the sequence's unprecedented length. As a result, methods that simply retain initial tokens without employing interpolation become ineffective in this regime, and we propose to address this limitation in future research.

\section{Related Works}

In addition to methods based on position interpolation, this section discusses related works of other approaches. 

\textbf{Retrieval-based} methods employ an external memory component to store extended historical context, alongside retrieval mechanisms to access pertinent documents during inference~\cite{longllama,wang2023augmenting,borgeaud2022improving}. Implementing these strategies often requires significant changes to the LLM architecture. By comparison, our approach introduces only minimal changes, specifically through slight adjustments to positional embeddings. Furthermore, our method extends beyond standard retrieval tasks, enabling it to address a wider range of long-context challenges such as lengthy document summarization and few-shot learning.

\textbf{Attention-based context window extensions}. In addition to techniques such as positional embedding interpolation, certain studies extend the input context using the standard LLM context window by adjusting attention mechanisms~\cite{han2023lminfinite, streamingllm,parallelwindow}. Central to these methods is the use of innovative attention masks to address the problem of attention explosion that arises from the introduction of new positions. These approaches are not mutually exclusive, but rather serve as complements to positional interpolation strategies.

\textbf{Fine-tuning based approaches} aim to optimize the adaptation of pre-trained LLMs to extended context lengths by adjusting position embeddings. For instance, methods such as Code LLaMA~\cite{codellama}, LLaMA2 Long~\cite{xiong2023effective}, and ScaledRoPE~\cite{liu2023scaling} employ exceptionally large base values for RoPE and proceed to fine-tune at the desired sequence length. In contrast, our approach provides adaptability across a range of target lengths, enabling support for lengths exceeding 2M tokens. With the increasing computational demands of fine-tuning for lengthy contexts (specifically above 128k), recent efforts like LongLoRA~\cite{longlora} and PoSE~\cite{zhu2023pose} aim to reduce the associated GPU requirements. Our technique is complementary to these efficiency-oriented fine-tuning strategies.

\section{Conclusion}

In this paper, we introduce LongRoPE, a technique that significantly increases the context window of LLMs to an unparalleled 2048k tokens, all while preserving their effectiveness within the original, shorter context ranges. By leveraging two types of non-uniform characteristics in RoPE positional embeddings—identified through an optimized evolutionary search process—we achieve two key advantages: strong initialization for subsequent fine-tuning, and an 8$\times$ expansion of the context window even in the absence of fine-tuning. Furthermore, we outline a staged extension approach that utilizes LLMs fine-tuned at a 256k context window, ultimately scaling up to 2048k context capability without requiring additional fine-tuning procedures. Our comprehensive experimental results demonstrate LongRoPE's robustness. We anticipate that our LongRoPE-2048k models will unlock numerous applications for long-context processing and stimulate continued research in this area.

\section*{Broader Impacts} The research described in this paper aims to contribute to the advancement of Machine Learning as a discipline. While our findings could lead to various societal effects, we do not identify any particular outcomes that require special attention at this time.

	\balance

\end{document}